This section provides an overall summary of the security risks
associated with the vulnerabilities confirmed during the assessment
of the OWASP Juice Shop application.

\subsection*{Overall Risk Assessment}

The assessment identified eight confirmed vulnerabilities across
authentication, authorization, input validation, client-side security,
request integrity, and redirect handling.

Three findings were classified as Critical, one as High, and four as
Medium. The Critical and High severity findings affect important
application security controls and should be prioritized for
remediation.

The risk assessment presented in this report is based on the
vulnerabilities that were identified and manually validated during
the engagement.

\subsection*{Severity Rating Definitions}

\begin{table}[H]
\centering
\renewcommand{\arraystretch}{1.3}
\begin{tabular}{|l|p{10cm}|}
\hline
\textbf{Severity} & \textbf{Definition} \\
\hline
Critical &
Vulnerabilities with the potential to cause severe compromise of
authentication, authorization, application functionality, or
sensitive data, depending on the affected component and exploitation
conditions. \\
\hline
High &
Vulnerabilities that can result in significant unauthorized access,
data exposure, account compromise, or execution of attacker-controlled
content under the applicable exploitation conditions. \\
\hline
Medium &
Vulnerabilities that weaken application security controls or enable
unauthorized actions or further attacks, but generally have more
limited impact or require additional conditions. \\
\hline
Low &
Vulnerabilities with limited direct security impact, such as minor
information disclosure or configuration weaknesses. \\
\hline
\end{tabular}
\caption{Severity rating definitions used in this report}
\end{table}

\subsection*{Risk Distribution}

\begin{table}[H]
\centering
\begin{tabular}{|l|l|l|}
\hline
\textbf{Severity} & \textbf{Number of Findings} & \textbf{Percentage} \\
\hline
Critical & 3 & 37.5\% \\
\hline
High & 1 & 12.5\% \\
\hline
Medium & 4 & 50\% \\
\hline
Low & 0 & 0\% \\
\hline
\textbf{Total} & \textbf{8} & \textbf{100\%} \\
\hline
\end{tabular}
\caption{Distribution of confirmed findings by severity}
\end{table}

\subsection*{Risk Matrix}

\begin{table}[H]
\centering
\renewcommand{\arraystretch}{1.3}
\begin{tabular}{|l|l|l|}
\hline
\textbf{Finding ID} & \textbf{Likelihood of Exploitation} &
\textbf{Business Impact} \\
\hline
SQLI-01 & High & Critical \\
\hline
BAC & High & Critical \\
\hline
AUTH & Medium & Critical \\
\hline
SQLI-02 & High & High \\
\hline
XSS & Medium & Medium \\
\hline
IDOR & High & Medium \\
\hline
OR & Medium & Medium \\
\hline
CSRF & Medium & Medium \\
\hline
\end{tabular}
\caption{Risk matrix showing likelihood and business impact for each
confirmed finding}
\end{table}

\subsection*{Key Risk Observations}

\begin{itemize}

    \item Three Critical-severity vulnerabilities were confirmed
    during the assessment. These findings affect authentication,
    authorization, and account recovery functionality.

    \item The BAC finding demonstrates a weakness in authorization
    controls that can allow unauthorized privilege-related actions
    through the affected registration functionality.

    \item SQLI-01 demonstrates that the login functionality can be
    bypassed through SQL Injection, resulting in unauthorized
    administrative access.

    \item The AUTH finding demonstrates that the password recovery
    mechanism relies on a predictable security question, allowing an
    attacker to take over the affected account without knowledge of
    the original password. If the same weakness affected a privileged
    account, it could potentially lead to unauthorized administrative
    access. This finding was classified as Critical severity based on
    its CVSS score and the potential for full account takeover.

    \item SQLI-02 was classified as High severity. It does not
    require authentication, but its confirmed impact was limited to
    unauthorized read access to database records rather than a
    full authentication bypass.

    \item The IDOR finding demonstrates that server-side authorization
    checks are insufficient for the affected basket resource, as the
    basket identifier can be manipulated to access another user's
    resource. This finding was classified as Medium severity based on
    the confidentiality-only impact and the requirement for an
    authenticated session.

    \item The XSS, CSRF, and Open Redirect findings demonstrate
    additional Medium-severity weaknesses affecting client-side
    execution, request integrity, and redirect handling respectively.

    \item No Low-severity vulnerabilities are included in the final
    confirmed findings because the assessment focused on findings
    that were successfully identified and manually validated with
    demonstrable security impact.

\end{itemize}

\subsection*{Conclusion of Risk Assessment}

The confirmed vulnerabilities demonstrate security weaknesses in
multiple application controls, including authentication,
authorization, input validation, client-side security, and request
integrity.

The Critical and High severity findings should be addressed before
the application is deployed in an environment where it processes
real user data or provides access to sensitive functionality.